\chapter{前提知識}
\label{chap:background}

本章では，獣害対策トレイルカメラシステムの設計と評価に必要な基礎事項を述べる．
まず，野生動物モニタリングにおけるトレイルカメラの役割と撮影の仕組みを説明する．
次に，画像認識による動物検出の基礎と，本研究で扱う検出モデルの概要を整理する．
さらに，撮影や推論を担うマイコンとシングルボードコンピュータ，および無線通信技術であるWi-Fi HaLowについて説明する．
個別の関連研究は第\ref{chap:related}章，本研究で採用する構成と設定は第\ref{chap:system}章，実験条件と結果は第\ref{chap:evaluation}章で述べる．

\section{トレイルカメラの概要}
\label{sec:background-trail-camera}
\subsection{トレイルカメラとは}
獣害対策では，動物がいつ，どこに出現したかを把握することが，見回りや対策箇所の検討に必要となる．
トレイルカメラは，監視対象の場所に設置し，センサの反応などを契機として自動的に画像を記録する装置である．
人が常時その場にいなくても画像を収集できるため，野生動物の出現状況や活動を調べる手段として用いられる\cite{Beery2019,IEICE2020}．
撮影日時と設置場所を画像に対応付けることで，画像に含まれる動物と，その出現時刻・地点を合わせて確認できる．

監視情報の利用には，蓄積した画像を事後的に分析する利用と，出現情報を速やかに伝えて対応を支援する利用がある．
後者では，撮影できることに加えて，画像の回収，判定，利用者への通知までの時間が重要となる．
画像の記録と通知は獣害対策を支援する機能であり，動物の検出そのものと，追い払いや侵入防止などによる被害の抑止とは区別して評価する必要がある．

\subsection{PIRセンサによる動体検知}
\label{sec:background-capture}
PIR（Passive Infrared）センサは，受光する赤外線の変化を利用する受動型のセンサである．
焦電型のPIRセンサでは，検知範囲内の熱源の移動などに伴う赤外線の変化を電気的な信号として取り出す．
検知は対象と背景からの赤外放射や対象の動きに影響され，対象が静止している場合などには反応しないことがある\cite{Welbourne2016,PanasonicPIR}．
この信号を撮影のトリガーとすることで，カメラを常時撮影させる構成とは異なる動作が可能になる．

Welbourneら\cite{Welbourne2016}は，PIRセンサによる検知を，物体表面からの赤外放射，フレネルレンズおよび焦電素子の関係から説明している．
フレネルレンズは複数の方向からの赤外線を素子に集め，空間的な検知領域を形成する．
対象の移動によって各素子に入射する放射が変化すると，撮影の契機となる信号が生じる．
したがって，PIRセンサは動物の体内温度と気温を直接比較する装置ではなく，検知しやすさは物体表面と背景の放射の違いや，検知領域を横切る動きに左右される．

ただし，PIRセンサは動物の形状や種を認識しているわけではない．
日射や周囲の温度変化など，監視対象の動物以外に由来する変化にも反応する場合がある\cite{PanasonicPIR}．
背景と異なる放射を示す草木の揺れや，カメラ自体の揺れによる検知領域の変化も，動物以外の撮影を生じさせる要因となる\cite{Welbourne2016}．
また，検知と撮影の間には時間差があるため，検知した対象が撮影時に画角内にいるとは限らない．
本論文では，撮影されたにもかかわらず動物が写っていない画像を空撮影画像と呼ぶ．
人や車両などが写った画像は，画像の内容としては空ではないため，必要に応じて別の分類として扱う．

カメラトラップの画像確認では，空撮影画像の選別が負担となることが指摘されている\cite{Beery2019}．
センサの反応をそのまま動物の出現として通知するのではなく，画像の内容を調べる処理を加えることが，確認作業や不要な通知を減らすための基本となる．

\subsection{昼間・夜間の撮影}
昼間と夜間では，撮影に利用できる光と得られる画像の特徴が異なる．
夜間の撮影では，赤外線照明を用いる構成などがあり，昼間の画像と同じ色や明るさの情報を得られるとは限らない．
ここで，撮影用の赤外線照明と，熱源の変化を捉えるPIRセンサとは役割が異なる．
前者はカメラで像を得るための照明であり，後者は撮影を開始する契機を与えるセンサである．

動物が画像内で占める大きさ，姿勢，草木による遮蔽，明るさやぶれは，目視確認と自動検出の双方で考慮すべき条件である．
Beeryらの画像解析では，夜間画像の確認や，枝・岩を動物と判定する誤検出の事例も示されている\cite{Beery2019}．
このため，昼夜を合わせた単一の評価値だけでなく，撮影条件を分けた評価によって，利用環境に応じた性能差を確認する必要がある．

\subsection{画像回収型と通信型}
画像回収型のトレイルカメラは，撮影した画像をSDカードなどの記憶媒体に保存し，利用者が現地で回収して確認する．
この方式では撮影地点に広域通信の接続環境を必要としない一方，画像を確認できる時点が回収の時期に依存する．
監視地点が増加すると，各地点への移動と画像確認の作業も運用上の負担となる．

通信型では，撮影画像や検出情報をネットワーク経由で遠隔地へ送る．
例えば，Saitoらは検出情報をサーバ経由で通知する構成を扱い\cite{IEICE2020}，Whytockらは衛星通信によって画像の判定ラベルとメタデータを通知する構成を報告している\cite{Whytock2023}．
通信型であっても，画像そのものを転送する方式と，判定結果だけを転送する方式とでは，利用者が確認できる情報と転送データ量が異なる．
したがって，通信機能の有無だけでなく，何を，どの時点で，どの経路に送るかを設計する必要がある．

\section{画像認識による動物検出}
\label{sec:background-detection}
\subsection{物体検出の概要}
画像分類は，画像に対して動物種などのカテゴリを割り当てる処理である．
これに対して物体検出は，画像中の対象の位置とカテゴリを推定する処理であり，対象を囲む矩形領域である検出枠と，クラス，信頼度スコアなどを出力する．
1枚の画像に複数の対象が存在する場合は，複数の検出結果が得られる．

動物の有無を調べることと，写っている動物の種を識別することも区別される．
Beeryらは，汎用的な動物検出器で動物の位置を求め，その領域を対象とする分類器と組み合わせる処理を示している\cite{Beery2019}．
空撮影画像を除く段階では動物の有無が主要な判定対象となるが，種別に通知や対応を変える場合には，種を識別するための追加の処理が必要となる．

検出モデルの出力から対象を採用するかどうかを決める際には，信頼度スコアに対する閾値を用いる．
画像$x$における対象クラスの検出候補の最大スコアを$s(x)$，判定閾値を$\tau$とすると，画像単位の判定は，例えば次のように表せる．
\begin{equation}
  \hat{y}(x)=
  \begin{cases}
    1 & \text{対象の検出候補が存在し，}s(x)\geq\tau,\\
    0 & \text{それ以外．}
  \end{cases}
  \label{eq:background-threshold}
\end{equation}
ここで，$\hat{y}=1$は対象あり，$\hat{y}=0$は対象なしを意味する．
信頼度は検出結果を選別するためのスコアであり，その値を実環境における正答確率と直ちに同一視することはできない．

同じ検出候補に対して閾値を下げると，採用される候補が増え，見逃しを抑えられる一方で誤検出を増やす可能性がある．
閾値を上げると逆に，低いスコアの対象を除外するため，動物が写った画像を残せない可能性がある．
前段の判定で転送対象から除外した画像は，後段の検出器だけでは再判定できない．
そのため，段階的な選別では，前段での見逃しと，後段へ渡す画像量の両方を評価することが重要となる．

\subsection{YOLO}
YOLO（You Only Look Once）は，画像全体から物体の位置とクラスを一つのニューラルネットワークで推定する物体検出手法として，Redmonら\cite{Redmon2016}が提案した．
原論文では，検出枠の座標とクラスの確率の推定を一体化し，画像に対する一度のネットワーク評価で複数の物体を検出する考え方を示している．
この考え方は，撮影画像を逐次処理する監視システムで検出器を選ぶ際の基礎となる．
ただし，原論文のネットワーク構造や処理速度を，後続のモデルの仕様や実行性能として用いることはできない．

本研究で扱う一次選別用の検出モデルは，YOLOv8nである．
YOLOv8は物体検出などに用いられるモデル群であり，YOLOv8nはその小規模なモデルである．
COCOデータセットで事前学習された物体検出モデルは，80クラスを対象としている\cite{UltralyticsYOLOv8}．
汎用データセットのクラスと，獣害対策で監視したい動物種とは必ずしも一致しないため，クラス名の一致だけで監視対象への検出性能を保証することはできない．
本研究では，種の厳密な識別と，動物が写っている可能性のある画像を残す処理を分けて考える．

\subsection{MegaDetector}
本研究では再判定用のモデルとしてMegaDetector v5aを扱う．
MegaDetectorは，カメラトラップ画像から動物，人，車両を検出するモデルであり，動物の種を識別するモデルとは異なる\cite{MegaDetectorDocs}．
動物の検出位置とスコアを利用することで，画像確認や空撮影画像の選別を支援できる\cite{Beery2019}．
モデルの実行時間やメモリ使用量は，モデル版，入力画像の大きさ，実行環境などに依存するため，名称のみからシステム上の処理性能を決めることはできない．
各モデルの配置と選択理由は第\ref{chap:system}章で述べ，性能は第\ref{chap:evaluation}章の条件下で評価する．

\section{マイコンとシングルボードコンピュータ}
\label{sec:background-computers}
\subsection{マイコンとシングルボードコンピュータの特徴}
マイクロコントローラ（マイコン）は，プロセッサ，メモリおよび周辺機器の制御機能を集積したデバイスであり，センサ入力の取得や機器の動作制御などに用いられる．
必要な処理をファームウェアとして実装し，機能ごとに動作と停止を制御することで，検知時のみ撮影・通信するような間欠動作を構成できる．
一方，シングルボードコンピュータは，プロセッサやメモリ，入出力機能などを一枚の基板にまとめた計算機である．
本研究で扱うRaspberry Pi 4のように，Linuxを基盤とするOS上で複数のプログラムを実行できる機器は，画像受信，ファイル管理，推論，通信を組み合わせる用途に利用できる\cite{RaspberryPiOS}．

両者の選択では，処理性能だけでなく，利用可能なメモリ，周辺機器との接続，起動時間，待機中の電力を考慮する．
マイコンでも推論は可能であり，シングルボードコンピュータでも省電力化は行えるため，名称だけで処理の可否や電力を決めることはできない．
野生動物監視では，必要な計算資源を各カメラに配置するか，共有する計算機に集約するかが，機器構成と電源設計に関わる．
動作時と待機時の電力を分け，実際に接続する周辺機器を含めて評価する必要がある\cite{Sato2022}．

\subsection{ESP32}
ESP32はEspressif Systemsが提供するマイコンの系列である．
本研究のカメラで扱うESP32-S3は，デュアルコアのプロセッサ，2.4\,GHz帯のIEEE 802.11b/g/n無線LAN，GPIOおよびカメラインタフェースなどを備える\cite{ESP32S3}．
GPIOはセンサの信号取得や周辺回路の制御に用いられ，カメラと無線通信を組み合わせることで，撮影した画像を外部へ転送する装置を構成できる．
チップ自体の機能と，それを搭載した基板のメモリ，電源回路，接続端子などの仕様は区別して扱う．

ESP32-S3には，通常動作に加えて，Light-sleepやDeep-sleepなどの動作モードが用意されている\cite{ESP32S3}．
待機中に停止する機能と，復帰に必要な機能を選ぶことで，撮影のない時間の消費電力を抑える設計が可能となる．
ただし，チップを低消費電力モードにしても，カメラ，記憶媒体，レギュレータなどへの給電が継続すれば，装置全体の待機電力はそれらの影響を受ける．
したがって，データシート上のチップの電流値と，完成したカメラの消費電流は分けて評価する．

\subsection{Raspberry Pi}
Raspberry Piには複数の製品系列があるが，本節ではシングルボードコンピュータであるRaspberry Pi 4 Model Bを対象とする．
同機器はプロセッサとメモリに加えて，USB，Ethernet，無線LAN，GPIOなどを備えている\cite{RaspberryPi4}．
Raspberry Pi OSはDebianを基盤とするOSであり，Linuxのファイル管理やネットワーク機能を利用できる\cite{RaspberryPiOS}．
この実行環境により，画像の受信サービスと推論プログラムを動作させ，判定結果や画像を保存・転送する処理を構築できる．

画像処理の実行時間は，機器のモデルだけでなく，入力画像，推論モデル，ソフトウェア構成などにも依存する．
また，常時受信を行う場合には，推論を実行していない時間にもOSや通信機能が動作する．
屋外での運用では，推論時の負荷に加えて待機時の電力と通信機器の電力も把握し，電源容量を検討する必要がある．
本研究で採用する基板，メモリ容量および実行環境は，第\ref{chap:system}章で具体的に示す．

\section{Wi-Fi HaLow}
\label{sec:background-halow}
\subsection{Wi-Fi HaLowの概要}
Wi-Fi HaLowは，IEEE 802.11ahに基づく無線LAN技術であり，1\,GHz未満の周波数帯を利用する．
IoT機器向けに通信範囲と省電力性を重視し，IPネットワーク上でデータを扱えることを特徴とする\cite{HaLowOverview,HaLowMM6108}．
監視システムへの適用では，カメラや集約地点を結ぶ無線区間の選択肢となる．

Tianら\cite{Tian2021}は，IEEE 802.11ahの物理層とMAC層の機能を整理し，多数の省電力端末を広い範囲で接続するための研究課題を述べている．
MAC層の機能には，端末群ごとに通信可能な時間を割り当てるRestricted Access Window（RAW）や，端末の起床時刻を調整するTarget Wake Time（TWT）がある．
RAWは多数端末の同時アクセスによる競合を抑えるために用いられ，TWTは通信のない時間の休止を支援する．
これらの機能を画像転送に適用する際には，端末数や画像発生頻度に応じた設定が必要となる．

Wi-Fi HaLowを利用するには，その方式に対応した無線機器が必要である．
ESP32-S3の内蔵無線LANは2.4\,GHz帯のIEEE 802.11b/g/nであり，それだけでIEEE 802.11ahの無線区間に接続できるわけではない\cite{ESP32S3,HaLowOverview}．
したがって，カメラに対応モジュールを接続する構成や，対応する中継機器を用いる構成など，無線区間と既存機器の接続方法を明確にする必要がある．

\subsection{従来のWi-Fiとの比較}
ここでは，従来のWi-Fiとして2.4\,GHz帯や5\,GHz帯を利用する無線LANを考える．
Wi-Fi HaLowは，それらより低い周波数帯を用いて広い範囲への接続を目指す技術である\cite{Tian2021,HaLowOverview}．
ただし，通信可能な距離は周波数だけでは決まらず，送信電力，アンテナ，受信感度，地形や障害物などに依存する．
そのため，ある条件下で報告された距離を，農地や山林での通信保証として用いることはできない．

比較の際には，通信距離に加えて，画像を送り終えるまでの時間，通信成功率，待機時と送信時の電力を評価する．
通信範囲を広げられても，転送に長い時間を要する場合には，通知遅延や1回の撮影に伴う電力量への影響が生じる．
省電力機能を備えることと，特定の撮影頻度・データ量で装置全体の電力量が減ることは区別する必要がある．
本研究で通信方式を比較する際には，同じ画像と設置条件を用いて，画像転送に必要な性能と電源条件の両面から適用性を検討する．
